\section{Construction of the expectation operator}

We wish to define the notion of the
$\textbf{expected value}$, or $\textbf{expectation}$,
of a random variable $X$, which will be denoted
$\mathbf{E} X$ (or $\mathbf{E}(X)$). In measure
theory this is denoted $\int X dP$ and is called the
``Lebesgue integral''. It is one of the most important
concepts in all of mathematical analysis! So time
invested in understanding it is time well-spent.

The idea is simple. For bounded random variables, we
want the expectation to satisfy three properties:
First, the expectation of an indicator variable $1_A$,
where $A$ is an event, should be equal to
$\mathrm{P}(A)$. Second, the expectation operator
should be linear i.e., should satisfy
$\mathbf{E}(aX+bY)=a\mathbf{E} X+b\mathbf{E} Y$ for
real numbers $a,b$ and r.v.'s $X,Y$. Third, it should
be monotone, i.e., if $X\le Y$ (meaning
$X(\omega)\le Y(\omega)$ for all $\omega\in \Omega$)
then $\mathbf{E} X\le \mathbf{E} Y$.

For unbounded random variables, we will also require
some kind of continuity, but let's treat the case of
bounded case first. It turns out that these properties
determine the expectation/Lebesgue integral operator
uniquely. Different textbooks may have some variation
in how they construct it, but the existence and
uniqueness are really the essential facts.

\begin{thm}
\label{thm-expecproperties}
Let $(\Omega, {\cal F}, \mathrm{P})$ be a probability
space. Let $B_\Omega$ denote the class of bounded
random variables. There exists a unique operator
$\mathbf{E}$ that takes a r.v. $X\in B_\Omega$ and
returns a number in $\R$, and satisfies:

\begin{enumerate}[label=\arabic*.]
\item If $A\in {\cal F}$ then
$\mathbf{E}(1_A) = \mathrm{P}(A)$.

\item If $X,Y\in B_\Omega$, $a,b\in\R$ then
$\mathbf{E}(aX+bY) = a\mathbf{E}(X)+b\mathbf{E}(Y)$.

\item If $X,Y\in B_\Omega$ and $X\ge Y$ then
$\mathbf{E}(X)\ge \mathbf{E}(Y)$.
\end{enumerate}

\end{thm}

\begin{proof}[Sketch of proof]
Call $X$ a $\textbf{simple}$ function if it is of the
form $X=\sum_{i=1}^n a_i 1_{B_i}$, where
$a_1,\ldots,a_n\in \R$ and $B_1,\ldots,B_n$ are
disjoint events. For such r.v.'s define
$\mathbf{E}(X)=\sum a_i \mathrm{P}(B_i)$. Show that
the linearity and monotonicity properties hold, and so
far uniqueness clearly holds since we had no choice in
how to define $\mathbf{E}(X)$ for such functions if
we wanted the properties above to hold. Now for a
general bounded r.v. $X$ with $|X|\le M$, for any
$\epsilon>0$ it is possible to approximate $X$ from
below and above by simple functions $Y \le X \le Z$
such that $\mathbf{E}(Z-Y) < \epsilon$. This suggests
defining

\begin{equation*} \tag{5}\label{eq:expec-def}
\mathbf{E}(X) = \sup \{ \mathbf{E}(Y) : Y\textrm{ is a simple function such that }Y \le X \}.
\end{equation*}

\noindent By approximation, the construction is shown to still
satisfy the properties in the Theorem and to be
unique, since $\mathbf{E}(X)$ is squeezed between
$\mathbf{E}(Y)$ and $\mathbf{E}(Z)$, and these can be
made arbitrarily close to each other.
\end{proof}

\noindent We can now extend the definition of the expectation
operator to non-negative random variables. In that
case we still define $\mathbf{E} X$ by eq.
\eqref{eq:expec-def}. This can be thought of as a kind of
``continuity from below'' axiom that is added to the
properties 1--3 above, although we shall see that it
can be reformulated in several equivalent ways. Note
that now $\mathbf{E} X$ may sometimes be infinite.

Finally, for a general random variable $X$, we
decompose $X$ as a difference of two non-negative
r.v.'s by writing

\[
X = X_+ - X_-,
\]

\noindent where $X_+ = \max(X,0)$ is called the
$\textbf{positive part of } X$ and $X_-=\max(-X,0)$
is called the $\textbf{negative part of } X$.

We say that $X$$\textbf{has an expectation}$ if the
two numbers $\mathbf{E} X_-, \mathbf{E} X_+$ are not
both $\infty$. In this case we define

\[
\mathbf{E} X = \mathbf{E} X_+ - \mathbf{E} X_-.
\]

\noindent This is a number in $\R\cup\{ -\infty,\infty\}$. If
both $\mathbf{E} X_-, \mathbf{E} X_+$ are $<\infty$,
or in other words if $\mathbf{E} |X|<\infty$ (since
$|X|=X_+ + X_-$), we say that $X$$\textbf{has finite expectation}$ or is
$\textbf{integrable}$.

\begin{thm}
Suppose $X,Y\ge 0$ or $X,Y\le 0$ or
$\mathbf{E} |X|, \mathbf{E}|Y|<\infty$. Then:

\begin{enumerate}[label=\arabic*.]
\item If $X$ is a simple function then the definition
\eqref{eq:expec-def} coincides with the original
definition, namely
$\mathbf{E}(\sum_i a_i 1_{B_i}) = \sum_i a_i
\mathrm{P}(B_i)$.

\item $\mathbf{E}(aX+bY+c)=a\mathbf{E} X + b\mathbf{E} Y+c$
for any real numbers $a,b,c$, where in the case where
$\mathbf{E}(X)=\mathbf{E}(Y) = \pm \infty$, we
require $a,b$ to have the same sign in order for the
right-hand side of this identity to be well-defined.

\item If $X\ge Y$ then $\mathbf{E} X\ge \mathbf{E} Y$.
\end{enumerate}

\end{thm}

\begin{proof}
See [Dur2010], section 1.4.
\end{proof}

\noindent $\textbf{Remark.}$ See \href{http://en.wikipedia.org/wiki/Lebesgue_integration\#Intuitive_interpretation}{the Wikipedia article on
Lebesgue
integration}
for a nice description of the difference in approaches
between the more familiar Riemann integral and the
Lebesgue integral.
